\scanpage{13}
\mspage{6}
% letterhead
\begin{center}
ACADEMIE D'AIX-MARSEILLE\\
\rule{4em}{0.4pt}\\
INSPECTION ACADEMIQUE\\
DE VAUCLUSE\\
\rule{1.6em}{0.4pt}
\end{center}
% ll. 1–9
mêmes objets que $[\widetilde\pi]\backslash B$ (où $[\widetilde\pi]$ est le foncteur composé de $\widetilde\pi:B\longrightarrow\mathrm{Ord}$ avec le foncteur oubli $\mathrm{Ord}\longrightarrow\mathrm{Ens}$), mais plus de flèches, une flèche de $(b,E)$ dans $(b',E')$ $(E\in\widetilde\pi(b),E'\in\widetilde\pi(b'))$ étant une flèche $u:b\longrightarrow b'$ telle que \deleted{$\unsure{u(E)|u_*E}$} $\widetilde\pi(u)(E)\subset E'$ (\unsure{ald|ald.} $\widetilde\pi(u)(E)=E'$).
% slanted
\margin{\uline{NB} $\widetilde\pi\backslash B$ est la cat. cofibrée sur $B$ définie par le foncteur composé $B\xrightarrow{\widetilde\pi}\mathrm{Ord}\longrightarrow\mathrm{Cat}$}
% ll. 10–12
Et \unsure{on a} un foncteur canonique
\[
 \begin{aligned}
 \widetilde\pi\backslash B&\longrightarrow(\underbrace{\pi\backslash B}_Z)^\wedge\\
 (b,\underset{\substack{\rotatebox{-90}{$\in$}\\\mathcal P^*(\pi(b))}}{E})&\longmapsto
 \widetilde E=\coprod_{a\in E}a\subset\coprod_{a\in\underset{\substack{\rotatebox{90}{$=$}\\\pi(b)}}{Z_b}}a .
 \end{aligned}\tag*{(*)}
\]
% ll. 13–19
Je voudrais des conditions sur $A$, ou encore sur $B=A^{\mathrm o}$, qui assurent que \inserted{quel que soit $\pi\in\operatorname{Ob}A^\wedge=\operatorname{Ob}B^\vee$,} le foncteur $(*)$ est pleinement fidèle$^*$, et son image essentielle est \deleted{\unsure{la plus|le plus}} stable par $\varprojlim$ finies. Par ailleurs, il est clair que cette image essentielle contient $Z$.

% ll. 20–21
Je m'aperçois qu'il est déraisonnable de s'attendre que le foncteur $(*)$ soit
% archive sheet number: 12
